\section{Specification deltas (documentation vs enforced code)}
\label{sec:deltas}

This section collects the places where the repository's own prose (the
top-level README, module and contract doc comments, deployment-script
comments, and the vendored-patch record) states a property that differs from
the one the code enforces. Each item gives (i) the documentation location,
(ii) the code location, and (iii) which of the two the client intends to be
authoritative, judged only from evidence inside the repository (the decision
log and migration notes under \code{.scratch/\allowbreak{}pq-shielded-rollup/\allowbreak{}}, in-code
rationale, and compile-time guards). Where the repository records no intent,
the item says so. Everywhere else in this document the \emph{code} is the
specification; the items below mark where a reader of the documentation alone
would form a different model. Several deltas were already noted locally in
earlier sections; this section consolidates them. The umbrella informational
finding is \finding{I-02}.

%----------------------------------------------------------------------
\subsection{Summary}
\label{subsec:deltas-summary}

{\small
\begin{longtable}{@{}>{\raggedright\arraybackslash}p{0.4cm}>{\raggedright\arraybackslash}p{2.6cm}>{\raggedright\arraybackslash}p{4.2cm}>{\raggedright\arraybackslash}p{4.9cm}>{\raggedright\arraybackslash}p{1.9cm}@{}}
\toprule
\# & Topic & Documentation states & Code enforces & Intended \\
\midrule
\endhead
1 & Shielded and commitment hashes
  & SHA3-256 for $\mathsf{pk}_d$ and $\mathsf{nf}$; Keccak-256 for
    $\mathsf{cm}$ and for the nullifier tree
  & Poseidon2 ($\Hp$, $\Cp$) for $\mathsf{pk}_d$, $\mathsf{nf}$,
    $\mathsf{cm}$, the note tree and the nullifier tree; documented
    $\sim$96-bit sponge margin
  & Code \\
2 & Caller of \code{applyBlock}
  & an ``operator set allowed to call \code{applyBlock}''
  & no caller restriction of any kind
  & Not recorded \\
3 & Empty nullifier root
  & computed at deploy: Keccak chain, depth 256
  & pinned Poseidon2 constant, depth 96; no hashing at deploy
  & Code \\
4 & Vendored \code{p3-recursion} patch record
  & one file (\code{pcs.rs}); ``\code{ZK} false$\to$true''; ``evaluation
    truncation''; a pristine upstream copy in-tree
  & five changed code files; \code{ZK} is a const generic defaulting to
    \code{true}; no truncation; no pristine copy
  & Code \\
5 & Security level
  & ``$\sim$128-bit''; ``the 128-bit target''
  & \code{SECURITY\_LEVEL = 96}; level, regime, folding, rate, grinding and
    LDE are environment-overridable in library code
  & Code (96) \\
6 & \code{WhirVerifier} wire-format header
  & STATEMENT ``not consumed here''; eq-points carried in PROOF
  & STATEMENT supplies the initial constraint's opening points; the
    eq-points blob is gone from the wire
  & Code \\
7 & Fees on settlement
  & ``roots + fees move''; the chain verifies ``fee accounting''
  & fee total decoded and emitted only; no payable entry point
  & Not recorded \\
\bottomrule
\end{longtable}
}

%----------------------------------------------------------------------
\subsection{Hash instantiation of the shielded and commitment layers}
\label{subsec:delta-hash}

\begin{specdelta}[SHA3-256 / Keccak-256 in prose, Poseidon2 in code]
\label{delta:hash}
\textbf{Documentation.} The README states that ``nullifiers are SHA3-256''
\src{README.md:11-12}, labels the wallet column ``SHA3-256 nullifiers''
\src{README.md:23}, defines the note commitment as
``\code{Keccak-256(DOMAIN || value || rho || psi || pk\_d)}'', calls the
nullifier map ``a \textbf{Keccak-256} map'' and the nullifiers ``SHA3-256 over
\code{(DOMAIN || sk\_d || rho)}'' \src{README.md:31-38}, and describes
\code{crates/\allowbreak{}pq-hash} as ``SHA3-256 note/nullifier hashing''
\src{README.md:60}. The \code{pq-hash} layering table assigns SHA3-256 to the
shielded layer and Keccak-256 to the nullifier tree
\src{crates/\allowbreak{}pq-hash/\allowbreak{}src/\allowbreak{}lib.rs:16-21}, and the crate preamble names
\code{Sha3\_256Shielded} as the shielded hasher
\src{crates/\allowbreak{}pq-hash/\allowbreak{}src/\allowbreak{}lib.rs:11-12}. The note methods describe the
commitment as Keccak-256 and the nullifier as SHA3-256
\src{crates/\allowbreak{}shielded/\allowbreak{}src/\allowbreak{}note.rs:104-106,118-120}. The prover's layering
diagrams describe the shielded layer as SHA3-256 computed in-circuit on
Keccak-f[1600] rows with the \code{0x06} pad byte
\src{crates/\allowbreak{}prover/\allowbreak{}src/\allowbreak{}config.rs:16-18}
\src{crates/\allowbreak{}prover/\allowbreak{}src/\allowbreak{}whir.rs:30-31}.

\textbf{Code.} With $\Hp$ the Poseidon2 byte hash and $\Cp$ the Poseidon2
two-to-one compression of Section~\ref{sec:primitives}, the enforced
relations are
\begin{align}
  \mathsf{pk}_d &= \Hp\big(\code{DOMAIN\_PK}\,\|\,\mathsf{sk}_d\big)
     & &\text{\src{crates/\allowbreak{}shielded/\allowbreak{}src/\allowbreak{}keys.rs:40-42}}\\
  \mathsf{nf} &= \Hp\big(\code{DOMAIN\_NULLIFIER}\,\|\,\mathsf{sk}_d\,\|\,\rho\big)
     & &\text{\src{crates/\allowbreak{}shielded/\allowbreak{}src/\allowbreak{}note.rs:125-126}}\\
  \mathsf{cm} &= \Hp\big(\code{DOMAIN\_NOTE}\,\|\,\mathrm{LE64}(v)\,\|\,\rho\,\|\,\psi\,\|\,\mathsf{pk}_d\big)
     & &\text{\src{crates/\allowbreak{}shielded/\allowbreak{}src/\allowbreak{}note.rs:108-111}}\\
  E_0 &= 0^{32},\qquad E_h = \Cp(E_{h-1},E_{h-1})\quad(1\le h\le 96)
     & &\text{\src{crates/\allowbreak{}shielded/\allowbreak{}src/\allowbreak{}nullifier\_tree.rs:129-134}}
\end{align}
where the shielded hasher is \code{Poseidon2Shielded}, whose
\code{hash\_to\_digest}(\textit{dom}, \textit{parts}) is the commitment byte
hash of $\textit{dom}\,\|\,\textit{parts}$ with no length prefixes
\src{crates/\allowbreak{}pq-hash/\allowbreak{}src/\allowbreak{}shielded.rs:98-104}, and \code{hash\_pair} of
\code{Poseidon2Commitment} is $\Cp$ \src{crates/\allowbreak{}pq-hash/\allowbreak{}src/\allowbreak{}poseidon.rs:240-242}.
These hashers are the ones instantiated at every production site: the node's
committed tree and nullifier map \src{crates/\allowbreak{}node/\allowbreak{}src/\allowbreak{}state.rs:48,50,110-111},
the wallet wasm \src{crates/\allowbreak{}wallet-wasm/\allowbreak{}src/\allowbreak{}lib.rs:347,520,533,536}, and the
native client prover \src{crates/\allowbreak{}prover/\allowbreak{}src/\allowbreak{}client.rs:166,204}. In-circuit,
$\mathsf{nf}$ and $\mathsf{cm}$ are Poseidon2 sponge evaluations
\src{crates/\allowbreak{}prover/\allowbreak{}src/\allowbreak{}transfer.rs:456,580}, and the transfer module states
that SHA3-256 and Keccak-f[1600] are absent from the circuit
\src{crates/\allowbreak{}prover/\allowbreak{}src/\allowbreak{}transfer.rs:30-43}. \code{Sha3\_256Shielded} and
\code{Keccak256Commitment} remain exported
\src{crates/\allowbreak{}pq-hash/\allowbreak{}src/\allowbreak{}lib.rs:55,62} but are referenced only from test code
(e.g.\ \src{crates/\allowbreak{}shielded/\allowbreak{}src/\allowbreak{}nullifier\_tree.rs:361-364}); the
\code{prover::sha3\_block} module is declared
\src{crates/\allowbreak{}prover/\allowbreak{}src/\allowbreak{}lib.rs:56} and referenced by no other module.
Keccak-256 remains the hash of the \emph{settlement} layer's Merkle tree and
Fiat--Shamir transcript only \src{crates/\allowbreak{}prover/\allowbreak{}src/\allowbreak{}config.rs:66-90}
\src{crates/\allowbreak{}prover/\allowbreak{}src/\allowbreak{}whir.rs:77}.

\textbf{Margin.} The shielded hasher's documentation records the collision
bound as the Poseidon2 sponge's ``$\sim$96-bit classical margin at this
width'', not SHA3's \src{crates/\allowbreak{}pq-hash/\allowbreak{}src/\allowbreak{}shielded.rs:82-85}, and the
nullifier-tree address is the low 96 bits of $\mathsf{nf}$
\src{crates/\allowbreak{}shielded/\allowbreak{}src/\allowbreak{}nullifier\_tree.rs:56-64}. The crate-level statement
of ``the same 128-bit collision security'' \src{crates/\allowbreak{}pq-hash/\allowbreak{}src/\allowbreak{}lib.rs:23-28}
compares Keccak-256 with SHA3-256 and no longer describes the shielded layer.

\textbf{Intended: code.} The migration record cites the directive ``no sha3,
only poseidon2'' and lists the SHA3/Keccak removals as consensus-critical
\src{.scratch/\allowbreak{}pq-shielded-rollup/\allowbreak{}verifier-redesign.md:3456-3478}; the code
records operator acceptance of the reduced margin
\src{crates/\allowbreak{}pq-hash/\allowbreak{}src/\allowbreak{}shielded.rs:85}
\src{crates/\allowbreak{}prover/\allowbreak{}src/\allowbreak{}nullifier\_gadget.rs:17-18}
\src{crates/\allowbreak{}prover/\allowbreak{}src/\allowbreak{}transfer.rs:41-43}. Related enforcement gaps are
catalogued as \finding{V-07} (in-circuit sponge) and \finding{M-09}
(nullifier address width).
\end{specdelta}

%----------------------------------------------------------------------
\subsection{Who may settle}
\label{subsec:delta-operator}

\begin{specdelta}[No operator set on \code{applyBlock}]
\label{delta:operator}
\textbf{Documentation.} The README's trust statement lists ``the operator set
allowed to call \code{applyBlock}'' among the things the chain trusts
\src{README.md:160-162}.

\textbf{Code.} \code{applyBlock} is declared \code{external} with no modifier
\src{contracts/\allowbreak{}src/\allowbreak{}ShieldedPool.sol:144}; the contract stores no operator,
role or allow-list (its state is \code{verifier}, \code{blockNumber},
\code{currentRoot}, \code{currentNullifierRoot}, \code{feeRecipient})
\src{contracts/\allowbreak{}src/\allowbreak{}ShieldedPool.sol:68-89}, and the only \code{msg.sender}
comparison in the contract guards \code{withdrawFees}
\src{contracts/\allowbreak{}src/\allowbreak{}ShieldedPool.sol:183}. The enforced relation is, for any
caller,
\begin{align}
  \code{applyBlock}(s,\pi)\ \text{succeeds} \iff\;
   & \code{verifier.verify}(s,\pi)=\mathsf{true} \notag\\
   & \wedge\; \mathsf{rootBefore}(s)=\code{currentRoot} \notag\\
   & \wedge\; \mathsf{nullifierBefore}(s)=\code{currentNullifierRoot}
\end{align}
\src{contracts/\allowbreak{}src/\allowbreak{}ShieldedPool.sol:145-157}
after which $(\code{currentRoot},\code{currentNullifierRoot}) \leftarrow
(\mathsf{rootAfter}(s),\mathsf{nullifierAfter}(s))$ and \code{blockNumber} is
incremented \src{contracts/\allowbreak{}src/\allowbreak{}ShieldedPool.sol:165-170}. The contract's own
trust-model text treats the prover and sequencer as untrusted and gates every
transition solely on a verified proof
\src{contracts/\allowbreak{}src/\allowbreak{}ShieldedPool.sol:51-55}, which matches the code; the node's
settlement sender submits from a configured unlocked account and registers no
operator \src{.scratch/\allowbreak{}pq-shielded-rollup/\allowbreak{}decisions.md:1342-1355}.

\textbf{Intended: not recorded.} No decision-log entry or code path introduces
an operator set; the README sentence is the only occurrence. This
specification takes the enforced relation above (permissionless settlement)
as normative. Consequences are catalogued in \finding{V-01}, \finding{H-01}
and \finding{M-01}.
\end{specdelta}

%----------------------------------------------------------------------
\subsection{Genesis nullifier root}
\label{subsec:delta-empty-root}

\begin{specdelta}[Keccak depth-256 computed root in prose, Poseidon2 depth-96 constant in code]
\label{delta:empty-root}
\textbf{Documentation.} The constructor documentation states that passing
\code{bytes32(0)} makes the contract compute ``the prover's empty-map root
itself (depth 256, empty leaf zero, \code{empty[h] = keccak(empty[h-1] ||
empty[h-1])} \ldots{} $\sim$256 keccak at deploy)''
\src{contracts/\allowbreak{}src/\allowbreak{}ShieldedPool.sol:114-118}. The deployment script says
\code{bytes32(0)} ``makes ShieldedPool compute the empty sparse-tree root
itself; the genesis file may pin one explicitly''
\src{contracts/\allowbreak{}script/\allowbreak{}Deploy.s.sol:42-43}. The nullifier-tree module
documentation describes full 256-bit addressing, 256 levels, and a path through
bits $0,\dots,255$ \src{crates/\allowbreak{}shielded/\allowbreak{}src/\allowbreak{}nullifier\_tree.rs:14-22,45-47}.

\textbf{Code.} The contract pins
\begin{align*}
  \code{EMPTY\_NULLIFIER\_ROOT} = {}&\code{0x83265b7a2c40cc6eca237b6016b76460}\\
                                    &\code{\phantom{0x}6335cd0c6ff0732470e20a484d5a7329}
\end{align*}
(one 32-byte word, printed in two halves)
\src{contracts/\allowbreak{}src/\allowbreak{}ShieldedPool.sol:85-86}, documented in the same file as
the Poseidon2 empty-subtree chain at \code{NULLIFIER\_TREE\_DEPTH = 96}
\src{contracts/\allowbreak{}src/\allowbreak{}ShieldedPool.sol:79-84}. The constructor performs no
hashing:
\begin{equation}
  \code{currentNullifierRoot} =
  \begin{cases}
    \code{EMPTY\_NULLIFIER\_ROOT} & \code{genesisNullifierRoot\_} = 0^{256}\\
    \code{genesisNullifierRoot\_} & \text{otherwise}
  \end{cases}
\end{equation}
\src{contracts/\allowbreak{}src/\allowbreak{}ShieldedPool.sol:134-135}
The deployment script passes \code{bytes32(0)} unconditionally
\src{contracts/\allowbreak{}script/\allowbreak{}Deploy.s.sol:44} and reads only
\code{.genesis\_root\_hex} from the genesis file
\src{contracts/\allowbreak{}script/\allowbreak{}Deploy.s.sol:32}; no nullifier-root field is consulted.
On the Rust side \code{NULLIFIER\_TREE\_DEPTH = 96}
\src{crates/\allowbreak{}shielded/\allowbreak{}src/\allowbreak{}nullifier\_tree.rs:64}, the empty-map root is $E_{96}$
of Delta~\ref{delta:hash}, and two nullifiers agreeing on the low 96 address
bits are treated as the same address
\src{crates/\allowbreak{}shielded/\allowbreak{}src/\allowbreak{}nullifier\_tree.rs:56-63,79-83}. The pinned constant
equals the \code{nullifier\_before\_hex} of the committed block vector
\src{contracts/\allowbreak{}test/\allowbreak{}vectors/\allowbreak{}block\_vectors.json:4}.

\textbf{Intended: code.} The migration record states ``nullifier tree:
Keccak $\to$ Poseidon2, depth 256 $\to$ 96'' and
``\code{\_emptyNullifierRoot()} keccak loop deleted;
\code{EMPTY\_NULLIFIER\_ROOT} constant \ldots{} pinned''
\src{.scratch/\allowbreak{}pq-shielded-rollup/\allowbreak{}verifier-redesign.md:3469-3477}. The
constructor comment and the deployment-script comment predate that change.
The 96-bit address width is the subject of \finding{M-09}.
\end{specdelta}

%----------------------------------------------------------------------
\subsection{Vendored recursion tree}
\label{subsec:delta-vendor}

\begin{specdelta}[\code{PATCHES.md} under-reports and mis-describes the vendored diff]
\label{delta:vendor}
\textbf{Documentation.} \code{vendor/\allowbreak{}p3-recursion/\allowbreak{}PATCHES.md} records the
upstream revision \code{8f9876efb26cff0f4b837555b478f043c4314240}
\src{vendor/\allowbreak{}p3-recursion/\allowbreak{}PATCHES.md:3-4}, states that a pristine copy is kept
at \code{.scratch/\allowbreak{}vendor/\allowbreak{}p3-recursion-pristine} from which the complete diff
is reproducible \src{vendor/\allowbreak{}p3-recursion/\allowbreak{}PATCHES.md:17-22}, and lists a single
patch: \code{recursion/\allowbreak{}src/\allowbreak{}pcs/\allowbreak{}whir/\allowbreak{}uni/\allowbreak{}pcs.rs}, ``\code{const ZK: bool =
false} $\to$ \code{true}; hiding-aware \code{commit}, \code{commit\_preprocessing},
\code{get\_opt\_randomization\_poly\_commitment}, and evaluation truncation''
\src{vendor/\allowbreak{}p3-recursion/\allowbreak{}PATCHES.md:35-37}. It cites decision-log entries
D-046 and D-051 for the rationale \src{vendor/\allowbreak{}p3-recursion/\allowbreak{}PATCHES.md:13,59}.

\textbf{Code.} A recursive comparison of the vendored tree with a clean
checkout of the upstream repository at the recorded revision (excluding
\code{.git} and \code{target}) yields exactly the following differences
besides \code{PATCHES.md} and the \code{.cargo-ok} marker:

{\small
\begin{longtable}{@{}>{\raggedright\arraybackslash}p{4.6cm}>{\raggedright\arraybackslash}p{10.4cm}@{}}
\toprule
File & Change \\
\midrule
\endhead
\code{recursion/\allowbreak{}Cargo.toml}
  & \code{rand} gains the \code{sys\_rng} feature and a \code{spin} dependency
    is added \src{vendor/\allowbreak{}p3-recursion/\allowbreak{}recursion/\allowbreak{}Cargo.toml:45-54}; the masking
    RNG is seeded from OS entropy
    \src{vendor/\allowbreak{}p3-recursion/\allowbreak{}recursion/\allowbreak{}src/\allowbreak{}pcs/\allowbreak{}whir/\allowbreak{}uni/\allowbreak{}pcs.rs:419-422}. \\
\code{recursion/\allowbreak{}src/\allowbreak{}backend/\allowbreak{}whir.rs}
  & upstream rejects inputs with \code{config.is\_zk() != 0} (``supports only
    non-ZK STARK inputs'') at two sites; both checks are removed, the
    surrounding code following directly at
    \src{vendor/\allowbreak{}p3-recursion/\allowbreak{}recursion/\allowbreak{}src/\allowbreak{}backend/\allowbreak{}whir.rs:562-565,643-647}. \\
\code{recursion/\allowbreak{}src/\allowbreak{}input\_contract/\allowbreak{}whir.rs}
  & test only: \code{whir\_config(8)} becomes \code{whir\_config(9)} to match
    the masked stacked arity
    \src{vendor/\allowbreak{}p3-recursion/\allowbreak{}recursion/\allowbreak{}src/\allowbreak{}input\_contract/\allowbreak{}whir.rs:1257-1260}. \\
\code{recursion/\allowbreak{}src/\allowbreak{}pcs/\allowbreak{}whir/\allowbreak{}uni/\allowbreak{}pcs.rs}
  & \code{ZK} becomes a const generic of \code{WhirUniPcs} with default
    \code{true} \src{vendor/\allowbreak{}p3-recursion/\allowbreak{}recursion/\allowbreak{}src/\allowbreak{}pcs/\allowbreak{}whir/\allowbreak{}uni/\allowbreak{}pcs.rs:322},
    forwarded by \code{const ZK: bool = ZK}
    \src{vendor/\allowbreak{}p3-recursion/\allowbreak{}recursion/\allowbreak{}src/\allowbreak{}pcs/\allowbreak{}whir/\allowbreak{}uni/\allowbreak{}pcs.rs:967} (upstream:
    \code{const ZK: bool = false}); hiding-aware commit/open paths are gated
    on it. \\
\code{recursion/\allowbreak{}src/\allowbreak{}pcs/\allowbreak{}whir/\allowbreak{}uni/\allowbreak{}recursive\_pcs.rs}
  & the \code{RecursivePcs} impl is generalised over \code{const ZK: bool}
    \src{vendor/\allowbreak{}p3-recursion/\allowbreak{}recursion/\allowbreak{}src/\allowbreak{}pcs/\allowbreak{}whir/\allowbreak{}uni/\allowbreak{}recursive\_pcs.rs:361-368}. \\
\bottomrule
\end{longtable}
}

No evaluation truncation exists: \code{evaluations\_on\_domain} asserts
$|\mathit{domain}| \ge$ the committed height and zero-pads the coefficients
before the coset DFT
\src{vendor/\allowbreak{}p3-recursion/\allowbreak{}recursion/\allowbreak{}src/\allowbreak{}pcs/\allowbreak{}whir/\allowbreak{}uni/\allowbreak{}pcs.rs:538-554}. The
directory \code{.scratch/\allowbreak{}vendor/\allowbreak{}} does not exist in the working tree. The
decision log has no D-046 or D-051 entry; D-046 is only mentioned in passing
\src{.scratch/\allowbreak{}pq-shielded-rollup/\allowbreak{}decisions.md:170,369}. The ``false$\to$true''
description matches only the \emph{default}: the workspace instantiates the
settlement PCS with \code{ZK = false}
\src{crates/\allowbreak{}prover/\allowbreak{}src/\allowbreak{}whir.rs:86}
\src{crates/\allowbreak{}prover/\allowbreak{}src/\allowbreak{}settlement\_replay.rs:87} and the inner (client-proof)
PCS with the default \src{crates/\allowbreak{}prover/\allowbreak{}src/\allowbreak{}whir\_recursion.rs:106-113}, and
pins both choices with compile-time assertions
\src{crates/\allowbreak{}prover/\allowbreak{}src/\allowbreak{}lib.rs:102-113}. The item-2 analysis in
\code{PATCHES.md} cites line numbers that still hold
\src{vendor/\allowbreak{}p3-recursion/\allowbreak{}recursion/\allowbreak{}src/\allowbreak{}pcs/\allowbreak{}whir/\allowbreak{}uni/\allowbreak{}recursive\_pcs.rs:384}
\src{vendor/\allowbreak{}p3-recursion/\allowbreak{}recursion/\allowbreak{}src/\allowbreak{}verifier/\allowbreak{}batch\_stark.rs:1528}.

\textbf{Intended: code.} The tree that builds is the vendored one, and the
non-ZK settlement layer is a recorded, deliberate choice
\src{crates/\allowbreak{}prover/\allowbreak{}src/\allowbreak{}lib.rs:92-101}; \code{PATCHES.md} was not updated with
the later changes. The removed upstream ZK guard and the hiding claims are
the subject of \finding{L-07} and \finding{M-07}; the non-ZK settlement
witness is the subject of \finding{M-10}.
\end{specdelta}

%----------------------------------------------------------------------
\subsection{Security level and configuration knobs}
\label{subsec:delta-security}

\begin{specdelta}[``128-bit'' in prose, 96-bit target in code, overridable at run time]
\label{delta:security}
\textbf{Documentation.} \code{config.rs} states that the degree-5 challenge
field gives ``$\sim$128-bit conjecturable security''
\src{crates/\allowbreak{}prover/\allowbreak{}src/\allowbreak{}config.rs:28-30,46-50}. The \code{whir.rs} module
header states that ``\code{SECURITY\_LEVEL} is set directly to the 128-bit
target'' \src{crates/\allowbreak{}prover/\allowbreak{}src/\allowbreak{}whir.rs:43-46}.

\textbf{Code.} The degree-5 FRI \code{Config} of \code{config.rs} is used only
by the test-gated \code{spike} module \src{crates/\allowbreak{}prover/\allowbreak{}src/\allowbreak{}lib.rs:64-65}
\src{crates/\allowbreak{}prover/\allowbreak{}src/\allowbreak{}spike.rs:18}; the deployed path takes only $\F$, the
Keccak Merkle MMCS and the Keccak challenger from it
\src{crates/\allowbreak{}prover/\allowbreak{}src/\allowbreak{}whir.rs:60,346-350}. The settlement challenge field is
the degree-4 binomial extension $\K$ \src{crates/\allowbreak{}prover/\allowbreak{}src/\allowbreak{}whir.rs:70}, and
\begin{equation}
  \code{SECURITY\_LEVEL} = 96
  \qquad\text{\src{crates/\allowbreak{}prover/\allowbreak{}src/\allowbreak{}whir.rs:120}}
\end{equation}
with an in-code derivation that 128 bits exceeds the structural ceiling on
this field \src{crates/\allowbreak{}prover/\allowbreak{}src/\allowbreak{}whir.rs:97-119}; the recursion layer is tied
to the same constant \src{crates/\allowbreak{}prover/\allowbreak{}src/\allowbreak{}whir\_recursion.rs:118-123}.

The following environment variables are read by library code (not under
\code{cfg(test)} or a feature) whenever a configuration is built:

{\small
\begin{longtable}{@{}>{\raggedright\arraybackslash}p{4.4cm}>{\raggedright\arraybackslash}p{5.4cm}>{\raggedright\arraybackslash}p{5.2cm}@{}}
\toprule
Variable & Effect & Source \\
\midrule
\endhead
\code{WHIR\_SECURITY\_LEVEL} & replaces the 96-bit target, settlement and inner;
  parsed as \code{usize}, no bound
  & \src{crates/\allowbreak{}prover/\allowbreak{}src/\allowbreak{}whir.rs:126-131,211}
    \src{crates/\allowbreak{}prover/\allowbreak{}src/\allowbreak{}whir\_recursion.rs:212} \\
\code{WHIR\_SOUNDNESS\_REGIME} & \code{udr}/\code{capacity}/Johnson (default)
  & \src{crates/\allowbreak{}prover/\allowbreak{}src/\allowbreak{}whir.rs:182-188} \\
\code{WHIR\_INNER\_SOUNDNESS\_REGIME} & same, inner layer
  & \src{crates/\allowbreak{}prover/\allowbreak{}src/\allowbreak{}whir\_recursion.rs:186-193} \\
\code{WHIR\_FOLDING\_FINAL} & settlement folding factor $\in\{4,8\}$
  & \src{crates/\allowbreak{}prover/\allowbreak{}src/\allowbreak{}whir.rs:147-153} \\
\code{WHIR\_INNER\_RATE} & inner starting rate $\in\{1,2,3\}$
  & \src{crates/\allowbreak{}prover/\allowbreak{}src/\allowbreak{}whir\_recursion.rs:314-318} \\
\code{WHIR\_POW\_FLOOR}, \code{WHIR\_INNER\_POW\_FLOOR} & grinding floor
  & \src{crates/\allowbreak{}prover/\allowbreak{}src/\allowbreak{}settlement\_replay.rs:151-154}
    \src{crates/\allowbreak{}prover/\allowbreak{}src/\allowbreak{}whir\_recursion.rs:345-348} \\
\code{WHIR\_FINAL\_LDE} & settlement LDE height
  & \src{crates/\allowbreak{}prover/\allowbreak{}src/\allowbreak{}composed\_export.rs:1255-1258} \\
\bottomrule
\end{longtable}
}

The node reaches these on its settlement path
(\code{settlement\_bundle} \src{crates/\allowbreak{}node/\allowbreak{}src/\allowbreak{}main.rs:453}
$\to$ \src{crates/\allowbreak{}prover/\allowbreak{}src/\allowbreak{}composed\_export.rs:1237,1259}
$\to$ \src{crates/\allowbreak{}prover/\allowbreak{}src/\allowbreak{}settlement\_replay.rs:180}
$\to$ \src{crates/\allowbreak{}prover/\allowbreak{}src/\allowbreak{}whir.rs:209-211}). The grinding-budget search
bound uses the constant rather than the overridden value
\src{crates/\allowbreak{}prover/\allowbreak{}src/\allowbreak{}whir.rs:276}. On-chain, the WHIR schedule is read from
the bundle's CONFIG section rather than fixed by the deployed contract
(Delta~\ref{delta:whirhdr}), so the deployed verifier does not itself fix the
level a settled proof attains \finding{V-01}.

\textbf{Intended: code (96 bits, fixed).} The 96-bit rationale is explicit
\src{crates/\allowbreak{}prover/\allowbreak{}src/\allowbreak{}whir.rs:97-120}, and the overrides are documented as
experiment-only (``shipping below 96 bits is a protocol decision, never a
default'') \src{crates/\allowbreak{}prover/\allowbreak{}src/\allowbreak{}whir.rs:122-125,178-181}, yet they are
compiled into the production path; see \finding{I-01}.
\end{specdelta}

%----------------------------------------------------------------------
\subsection{Settlement-verifier header}
\label{subsec:delta-whirhdr}

\begin{specdelta}[STATEMENT section described as unconsumed]
\label{delta:whirhdr}
\textbf{Documentation.} The \code{WhirVerifier} header describes the bundle
as CONFIG, PROOF and ``STATEMENT (audit surface, not consumed here)''
\src{contracts/\allowbreak{}src/\allowbreak{}verifier/\allowbreak{}WhirVerifier.sol:15-17}; lists among PROOF
contents ``the zeta-derived eq\_points (D-072 phase 1)''
\src{contracts/\allowbreak{}src/\allowbreak{}verifier/\allowbreak{}WhirVerifier.sol:26-29}; states that every domain
point is ``recomputed here, never read from the proof''
\src{contracts/\allowbreak{}src/\allowbreak{}verifier/\allowbreak{}WhirVerifier.sol:31-35}; and states that CONFIG is
deploy-time data that ``a deployment should pin''
\src{contracts/\allowbreak{}src/\allowbreak{}verifier/\allowbreak{}WhirVerifier.sol:22-25}. An inline comment repeats
``zeta-derived, proof-supplied, D-072 phase 1''
\src{contracts/\allowbreak{}src/\allowbreak{}verifier/\allowbreak{}WhirVerifier.sol:494-498}.

\textbf{Code.} \code{verify} locates the STATEMENT section from its u32
little-endian word count after PROOF and bounds-checks it
\src{contracts/\allowbreak{}src/\allowbreak{}verifier/\allowbreak{}WhirVerifier.sol:244-257}; every opening round
receives it \src{contracts/\allowbreak{}src/\allowbreak{}verifier/\allowbreak{}WhirVerifier.sol:465-470} and sets the
initial WHIR constraint to mode~2 over that slice
\src{contracts/\allowbreak{}src/\allowbreak{}verifier/\allowbreak{}WhirVerifier.sol:503-513}; the raw slice is copied
into the \code{TerminalWeight} call frame
\src{contracts/\allowbreak{}src/\allowbreak{}verifier/\allowbreak{}WhirVerifier.sol:826-840}, which derives the
initial constraint's equality groups from the opening points it contains
\src{contracts/\allowbreak{}src/\allowbreak{}verifier/\allowbreak{}TerminalWeight.sol:38-40,108-125}. The per-round
eq-points blob is no longer on the wire
\src{contracts/\allowbreak{}src/\allowbreak{}verifier/\allowbreak{}WhirVerifier.sol:1101-1102}. Hence STATEMENT is
consumed: it supplies the opening points of each round's initial constraint.
The deployed verifier holds no CONFIG digest (its only pin is the satellite
code hash \src{contracts/\allowbreak{}src/\allowbreak{}verifier/\allowbreak{}WhirVerifier.sol:122,131,895}); a
CONFIG-pinning wrapper exists as \code{WhirVerifierV6}
\src{contracts/\allowbreak{}src/\allowbreak{}verifier/\allowbreak{}WhirVerifierV6.sol:44,63}, but the deployment
script constructs \code{WhirVerifier} and passes it to the pool
\src{contracts/\allowbreak{}script/\allowbreak{}Deploy.s.sol:41,45}.

\textbf{Intended: code.} The D-086 step~C comments at the consuming sites
describe the current behaviour
\src{contracts/\allowbreak{}src/\allowbreak{}verifier/\allowbreak{}WhirVerifier.sol:503-506,827-830,1101-1102}; the
header predates them. The opening points so supplied are not related to the
transcript's out-of-domain challenge, which is \finding{V-02}; the unpinned
CONFIG is \finding{V-01}.
\end{specdelta}

%----------------------------------------------------------------------
\subsection{Fees}
\label{subsec:delta-fees}

\begin{specdelta}[Fees described as moving on settlement]
\label{delta:fees}
\textbf{Documentation.} The README diagram says ``roots + fees move'' on
\code{applyBlock} \src{README.md:23-26}, and the security model lists ``fee
accounting'' among what the chain verifies \src{README.md:157-159}.

\textbf{Code.} \code{ShieldedPool} declares no \code{payable} function,
\code{receive} or \code{fallback} \src{contracts/\allowbreak{}src/\allowbreak{}ShieldedPool.sol:63-190}.
The block fee is decoded as $\sum_i \mathsf{fee}_i$ over the per-transfer fee
limbs \src{contracts/\allowbreak{}src/\allowbreak{}BlockStatement.sol:104-107} and is only emitted in
\code{BlockApplied} \src{contracts/\allowbreak{}src/\allowbreak{}ShieldedPool.sol:171-178};
\code{withdrawFees} forwards whatever balance the contract holds
\src{contracts/\allowbreak{}src/\allowbreak{}ShieldedPool.sol:182-189}. The contract's own documentation
says it does not check fees \src{contracts/\allowbreak{}src/\allowbreak{}ShieldedPool.sol:59-62}.

\textbf{Intended: not recorded.} See \finding{I-04}.
\end{specdelta}

%----------------------------------------------------------------------
\subsection{Pointer to the findings}

The deltas above are documentation defects in their own right, but several of
them mark exactly the places where the enforced relation is weaker than the
documented one. The security consequences of these deltas, and of the
enforcement gaps noted throughout Sections~\ref{sec:overview}--\ref{sec:application},
are catalogued in the audit report under \code{zk-findings/\allowbreak{}reports/\allowbreak{}}: findings
\finding{V-01}--\finding{V-08} (Critical), \finding{H-01}--\finding{H-03}
(High), and \finding{M-01}--\finding{M-11} (Medium), with the low-severity and
informational items (including \finding{I-01}, \finding{I-02} and
\finding{I-04} cited here) listed alongside them. This specification states
what the code enforces; the report states what that permits.
